\documentclass[../../main.tex]{subfiles}

\begin{document}

\section{Nikodym 度量空间: Vitali-Hahn-Saks 定理}\label{section:ch18-nikodym}

令 $(X,\mathcal{M},\mu)$ 为有限测度空间. 回忆两个可测集 $A$ 与 $B$ 的对称差 $A\Delta B$ 定义为
\begin{align*}
A\Delta B=[A\sim B]\cup[B\sim A]=[A\cup B]\sim[A\cap B].
\end{align*}
我们把以下集合等式的证明留作练习:
\begin{align}
(A\Delta B)\Delta(B\Delta C)=A\Delta C.\label{align:ch18-triangle}
\end{align}
通过对满足 $\mu(A\Delta B)=0$ 的 $A$, $B$ 定义 $A\simeq B$, 我们引入 $\mathcal{M}$ 上的关系 $\simeq$. 以上等式蕴涵着该关系是传递的, 而它显然是自反与对称的. 因此, 等价关系 $\simeq$ 诱导了 $\mathcal{M}$ 的一个等价类分解. 对 $A$, 记 $A$ 的等价类为 $[A]$. 在 $\mathcal{M}/\simeq$ 上定义 Nikodym 度量 $\rho_{\mu}$ 为
\begin{align*}
\text{对所有 }A,B\in\mathcal{M},\quad\rho_{\mu}([A],[B])=\mu(A\Delta B).
\end{align*}
我们从等式\eqref{align:ch18-triangle}推出 $\rho_{\mu}$ 是恰当定义的且三角不等式成立: 度量的其余两个性质是显然的. 我们称 $(\mathcal{M}/\simeq,\rho_{\mu})$ 为与测度空间 $(X,\mathcal{M},\mu)$ 相联系的 Nikodym 度量空间. 现在令 $\nu$ 为 $\mathcal{M}$ 上的有限测度. 它关于 $\mu$ 绝对连续. 对满足 $A\simeq B$ 的 $A$, $B\in\mathcal{M}$, 由于 $\mu(A\Delta B)=0$, $\nu(A\Delta B)=0$, 因此
\begin{align*}
\nu(A)-\nu(B)=[\nu(A\cap B)+\nu(A\sim B)]-[\nu(A\cap B)+\nu(B\sim A)]=\nu(A\sim B)-\nu(B\sim A)=0.
\end{align*}
于是我们定义 $\mathcal{M}/\simeq$ 上的 $\nu$ 为
\begin{align*}
\text{对所有 }A\in\mathcal{M},\quad\nu([A])=\nu(A).
\end{align*}
如同我们对 $L^{p}$ 空间所做的那样, 为简明起见, 我们将集合与等价类等同起来, 将 $(\mathcal{M}/\simeq,\rho_{\mu})$ 记为 $(\mathcal{M},\rho_{\mu})$.

Baire 范畴定理的一个推论(第 10 章的定理 7)告诉我们, 若完备度量空间上的实值函数序列逐点收敛于某个实值函数, 则存在该空间的一个点使得序列在该点是等度连续的. 为了在绝对连续测度序列的研究中使用这个结果, 我们现在证明 $(\mathcal{M},\rho_{\mu})$ 是完备度量空间且 $\mathcal{M}$ 上关于 $\mu$ 绝对连续的测度诱导了 $(\mathcal{M},\rho_{\mu})$ 上的一致连续函数.

\begin{theorem}\label{theorem:ch18-nikodym-complete}
令 $(X,\mathcal{M},\mu)$ 为有限测度空间. 则 Nikodym 度量空间 $(\mathcal{M},\rho_{\mu})$ 是完备的, 即每个 Cauchy 序列收敛.
\end{theorem}

\begin{proof}
观察到对 $A,B\in\mathcal{M}$,
\begin{align}
\mu(A\Delta B)=\int_{X}|\chi_{A}-\chi_{B}|\,d\mu.\label{align:ch18-char-distance}
\end{align}
定义算子 $T:\mathcal{M}\to L^{1}(X,\mu)$ 为 $T(E)=\chi_{E}$. 则\eqref{align:ch18-char-distance}断言算子 $T$ 是等距的, 即
\begin{align}
\text{对所有 }A,B\in\mathcal{M},\quad\rho_{\mu}(A,B)=\|T(A)-T(B)\|_{1}.\label{align:ch18-isometry}
\end{align}
令 $\{A_{n}\}$ 为 $(\mathcal{M},\rho_{\mu})$ 中的 Cauchy 序列. 则 $\{T(A_{n})\}$ 是 $L^{1}(X,\mu)$ 中的 Cauchy 序列. Riesz-Fischer 定理告诉我们, 存在函数 $f\in L^{1}(X,\mu)$ 使得在 $L^{1}(X,\mu)$ 中 $T(A_{n})\to f$ 且 $\{T(A_{n})\}$ 的子序列在 $X$ 上几乎处处逐点收敛到 $f$. 由于每个 $T(A_{n})$ 取值 $0$ 和 $1$, 若我们定义 $A_{0}$ 为 $X$ 中那些使得逐点收敛的子序列收敛到 $1$ 的点, 则在 $X$ 上几乎处处 $f=\chi_{A_{0}}$. 因此, 根据\eqref{align:ch18-char-distance}, 在 $(\mathcal{M},\rho_{\mu})$ 中 $A_{n}\to A_{0}$.
\end{proof}

\begin{lemma}\label{lemma:ch18-absolute-continuity-uniform}
令 $(X,\mathcal{M},\mu)$ 为有限测度空间而 $\nu$ 是 $\mathcal{M}$ 上的有限测度. 令 $E_{0}$ 为可测集而 $\varepsilon>0$ 与 $\delta>0$, 使得对任何可测集 $E$,
\begin{align}
\text{若 }\rho_{\mu}(E,E_{0})<\delta,\quad\text{则 }|\nu(E)-\nu(E_{0})|<\varepsilon/4.\label{align:ch18-lemma-proof1}
\end{align}
则对任何可测集 $A$ 和 $B$,
\begin{align}
\text{若 }\rho_{\mu}(A,B)<\delta,\quad\text{则 }|\nu(A)-\nu(B)|<\varepsilon.\label{align:ch18-lemma-proof2}
\end{align}
\end{lemma}

\begin{proof}
我们首先证明:
\begin{align}
\text{若 }\rho_{\mu}(A,\varnothing)<\delta,\quad\text{则 }\nu(A)<\varepsilon/2.\label{align:ch18-lemma-proof3}
\end{align}
观察到若 $D\subseteq C$, 则 $C\Delta D=C\sim D$. 令 $A$ 属于 $\mathcal{M}$ 且 $\rho_{\mu}(A,\varnothing)=\mu(A)<\delta$. 观察到,
\begin{align*}
[E_{0}\sim A]\Delta E_{0}=E_{0}\sim[E_{0}\sim A]=E_{0}\cap A\subseteq A.
\end{align*}
因此 $\rho_{\mu}(E_{0}\sim A,E_{0})=\mu([E_{0}\sim A]\Delta E_{0})\leqslant\mu(A)<\delta$, 根据假设\eqref{align:ch18-lemma-proof1},
\begin{align*}
\nu(E_{0})-\nu[E_{0}\sim A]<\varepsilon/4.
\end{align*}
我们从 $\nu$ 的分割性质推出
\begin{align*}
\nu(A\cap E_{0})=\nu(E_{0})-\nu(E_{0}\sim A)<\varepsilon/4.
\end{align*}
现在观察到
\begin{align*}
E_{0}\Delta[E_{0}\cup[A\sim E_{0}]]=[E_{0}\cup[A\sim E_{0}]]\sim E_{0}=A\sim E_{0}\subseteq A.
\end{align*}
因此, 用上面的方法,
\begin{align*}
\nu(A\sim E_{0})=\nu(E_{0}\cup[A\sim E_{0}])-\nu(E_{0})<\varepsilon/4.
\end{align*}
因此,
\begin{align*}
\nu(A)=\nu(A\cap E_{0})+\nu(A\sim E_{0})<\varepsilon/2.
\end{align*}
但对任何两个可测集, 由于 $\nu$ 是实值的与有限可加的,
\begin{align*}
\nu(A)-\nu(B)=[\nu(A\sim B)+\nu(A\cap B)]-[\nu(B\sim A)+\nu(A\cap B)]=\nu(A\sim B)-\nu(B\sim A).
\end{align*}
因此\eqref{align:ch18-lemma-proof3}蕴涵\eqref{align:ch18-lemma-proof2}.
\end{proof}

命题\cref{proposition:ch18-abs-cont-eps-delta}告诉我们 $\mathcal{M}$ 上的有限测度 $\nu$ 关于 $\mu$ 绝对连续当且仅当对每个 $\varepsilon>0$, 存在 $\delta>0$ 使得若 $\mu(E)<\delta$, 则 $\nu(E)<\varepsilon$. 这意味着若 $\nu$ 是有限的, 则 $\nu$ 关于 $\mu$ 绝对连续当且仅当集函数 $\nu$ 关于 Nikodym 度量在 $\varnothing$ 处连续. 然而, 我们从前一个引理推出, 若 $\mathcal{M}$ 上的有限测度 $\nu$ 关于 Nikodym 度量在 $\mathcal{M}$ 中的某个集合 $E_{0}$ 处连续, 则它在 $\mathcal{M}$ 上是一致连续的. 我们因此证明了以下命题.

\begin{proposition}\label{proposition:ch18-uniform-continuity}
令 $(X,\mathcal{M},\mu)$ 为有限测度空间, 而 $\nu$ 是 $\mathcal{M}$ 上关于 $\mu$ 绝对连续的有限测度. 则 $\nu$ 诱导了在与 $(X,\mathcal{M},\mu)$ 相联系的 Nikodym 度量空间上恰当定义的一致连续函数.
\end{proposition}

\begin{definition}[集合态收敛]\label{definition:ch18-setwise-convergence}
令 $(X,\mathcal{M})$ 为可测空间. $\mathcal{M}$ 上的测度序列 $\{\nu_{n}\}$ 称为在 $\mathcal{M}$ 上集合态地收敛到集函数 $\nu$, 若
\begin{align*}
\text{对所有 }E\in\mathcal{M},\quad\nu(E)=\lim_{n\to\infty}\nu_{n}(E).
\end{align*}
\end{definition}

\begin{definition}[一致绝对连续]\label{definition:ch18-uniform-abs-cont}
令 $(X,\mathcal{M},\mu)$ 为有限测度空间. $\mathcal{M}$ 上的有限测度序列 $\{\nu_{n}\}$ 称为关于 $\mu$ 一致绝对连续, 其中每个关于 $\mu$ 绝对连续, 若对每个 $\varepsilon>0$, 存在 $\delta>0$, 使得对任何可测集 $E$ 和任何自然数 $n$,
\begin{align*}
\text{若 }\mu(E)<\delta,\quad\text{则 }\nu_{n}(E)<\varepsilon.
\end{align*}
\end{definition}

不难看出, 用引理\cref{lemma:ch18-absolute-continuity-uniform}, $\mathcal{M}$ 的有限测度序列 $\{\nu_{n}\}$ 关于 $\mu$ 一致绝对连续当且仅当函数序列 $\{\nu_{n}:\mathcal{M}\to\mathbf{R}\}$ 关于 Nikodym 度量 $\rho_{\mu}$ 是等度连续的. 此外, 对每个自然数 $n$, Radon-Nikodym 定理告诉我们, 存在非负可积函数 $f_{n}$, $\mu$ 关于 $\nu_{n}$ 的 Radon-Nikodym 导数, 使得
\begin{align*}
\text{对所有 }E\in\mathcal{M},\quad\nu_{n}(E)=\int_{E}f_{n}\,d\mu.
\end{align*}
显然函数序列 $\{f_{n}\}$ 在 $X$ 上关于 $\mu$ 一致可积当且仅当测度序列 $\{\nu_{n}\}$ 关于 $\mu$ 一致绝对连续. 我们因此有以下命题.

\begin{proposition}\label{proposition:ch18-uniform-integrability-equivalence}
令 $(X,\mathcal{M},\mu)$ 为有限测度空间, 而 $\{\nu_{n}\}$ 是 $\mathcal{M}$ 的有限测度序列, 其中每个关于 $\mu$ 绝对连续. 则以下命题等价:
\begin{description}
\item[（i）] 测度序列 $\{\nu_{n}\}$ 关于测度 $\mu$ 一致绝对连续.
\item[（ii）] 函数序列 $\{\nu_{n}:\mathcal{M}\to\mathbf{R}\}$ 关于 Nikodym 度量 $\rho_{\mu}$ 等度连续.
\item[（iii）] Radon-Nikodym 导数序列 $\left\{\dfrac{d\mu}{d\nu_{n}}\right\}$ 在 $X$ 上关于测度 $\mu$ 是一致可积的.
\end{description}
\end{proposition}

\begin{theorem}\label{theorem:ch18-uniform-limit-measure}
令 $(X,\mathcal{M},\mu)$ 为有限测度空间, 而 $\{\nu_{n}\}$ 是 $\mathcal{M}$ 上关于 $\mu$ 一致绝对连续的有限测度序列. 若 $\{\nu_{n}\}$ 在 $\mathcal{M}$ 上集合态地收敛到 $\nu$, 则 $\nu$ 是 $\mathcal{M}$ 的关于 $\mu$ 绝对连续的测度.
\end{theorem}

\begin{proof}
显然, $\nu$ 是非负集函数. 有限可加集函数的集合态极限是有限可加的. 因此 $\nu$ 是有限可加的. 我们必须证明它是可数可加的. 令 $\{E_{k}\}_{k=1}^{\infty}$ 为可测集的不交族. 我们必须证明:
\begin{align}
\nu\left(\bigcup_{k=1}^{\infty}E_{k}\right)=\sum_{k=1}^{\infty}\nu(E_{k}).\label{align:ch18-uniform-limit-proof1}
\end{align}
若存在 $k$ 使得 $\nu(E_{k})=\infty$, 则根据 $\nu$ 的单调性, \eqref{align:ch18-uniform-limit-proof1}成立, 由于两边都是无穷的. 我们因此假设对所有 $k$, $\nu(E_{k})<\infty$. 根据 $\nu$ 的有限可加性,
\begin{align}
\text{对每个自然数 }n,\quad\nu\left(\bigcup_{k=1}^{\infty}E_{k}\right)=\sum_{k=1}^{n}\nu(E_{k})+\nu\left(\bigcup_{k=n+1}^{\infty}E_{k}\right).\label{align:ch18-uniform-limit-proof2}
\end{align}
令 $\varepsilon>0$. 根据序列 $\{\nu_{n}\}$ 关于 $\mu$ 的一致绝对连续性, 存在 $\delta>0$ 使得对可测集 $E$ 与任何自然数 $n$,
\begin{align}
\text{若 }\mu(E)<\delta,\quad\text{则 }\nu_{n}(E)<\varepsilon/2.\label{align:ch18-uniform-limit-proof3}
\end{align}
因此, 若 $\mu(E)<\delta$, 则 $\nu(E)<\varepsilon$. 由于 $\mu(X)<\infty$ 而 $\mu$ 是可数可加的, 存在自然数 $N$ 使得
\begin{align*}
\mu\left(\bigcup_{k=N+1}^{\infty}E_{k}\right)<\delta.
\end{align*}
根据 $\delta$ 的这个选择、\eqref{align:ch18-uniform-limit-proof2}以及每个 $\nu(E_{k})$ 的有限性我们得出:
\begin{align*}
\nu\left(\bigcup_{k=1}^{\infty}E_{k}\right)-\sum_{k=1}^{\infty}\nu(E_{k})<\varepsilon.
\end{align*}
因此\eqref{align:ch18-uniform-limit-proof1}得证. 因此 $\nu$ 是测度, 且我们从\eqref{align:ch18-uniform-limit-proof3}推出, 若 $\mu(E)=0$, 则 $\nu(E)=0$. 因此 $\nu$ 关于 $\mu$ 是绝对连续的.
\end{proof}

以下非凡的定理告诉我们, 在前一个定理的叙述中, 若序列 $\{\nu_{n}(E)\}$ 是有界的, 我们可以不用假设序列是一致绝对连续的: 一致绝对连续性是集合态收敛的结果. 该定理的证明是在一致有界原理、开映射定理之后 Baire 范畴定理的杰出成果.

\begin{theorem}[Vitali-Hahn-Saks]\label{theorem:ch18-vitali-hahn-saks}
令 $(X,\mathcal{M},\mu)$ 为有限测度空间, 而 $\{\nu_{n}\}$ 是 $\mathcal{M}$ 上的有限测度序列, 每个关于 $\mu$ 绝对连续. 假定 $\{\nu_{n}(X)\}$ 是有界的且 $\{\nu_{n}\}$ 在 $\mathcal{M}$ 上集合态地收敛于 $\nu$. 则序列 $\{\nu_{n}\}$ 关于 $\mu$ 一致绝对连续. 此外, $\nu$ 是 $\mathcal{M}$ 上关于 $\mu$ 绝对连续的有限测度.
\end{theorem}

\begin{proof}
根据定理\cref{theorem:ch18-nikodym-complete}, Nikodym 度量空间是完备的, 而 $\{\nu_{n}\}$ 诱导了在该度量空间上逐点收敛(即集合态)于函数 $\nu$ 的连续函数序列. $\nu$ 是实值的, 由于 $\{\nu_{n}(X)\}$ 是有界的. 我们从第 10 章的定理 7(Baire 范畴定理的一个推论)推出, 存在集合 $E_{0}\in\mathcal{M}$ 使得函数序列 $\{\nu_{n}:\mathcal{M}\to\mathbf{R}\}$ 在 $E_{0}$ 等度连续, 即对每个 $\varepsilon>0$, 存在 $\delta>0$, 使得对每个可测子集 $E$ 与每个自然数 $n$,
\begin{align*}
\text{若 }\rho_{\mu}(E,E_{0})<\delta,\quad\text{则 }|\nu_{n}(E)-\nu_{n}(E_{0})|<\varepsilon.
\end{align*}
由于这对每个 $\varepsilon>0$ 成立且每个 $\nu_{n}$ 是有限的, 我们从引理\cref{lemma:ch18-absolute-continuity-uniform}推出对每个 $\varepsilon>0$, 存在 $\delta>0$ 使得对每个可测子集 $E$ 与每个自然数 $n$,
\begin{align*}
\text{若 }\rho_{\mu}(E)<\delta,\quad\text{则 }\nu_{n}(E)<\varepsilon.
\end{align*}
因此 $\{\nu_{n}\}$ 是一致绝对连续的. 根据前一个定理, $\nu$ 是关于 $\mu$ 绝对连续的有限测度.
\end{proof}

\begin{remark}
当然, sigma 代数不是线性空间而测度不是线性算子. 虽然如此, 但是 Vitali-Hahn-Saks 定理与连续线性算子序列的逐点极限的连续性却惊人相似, 而 Baire 范畴定理是证明这些结果的基础. 我们也观察到引理\cref{lemma:ch18-absolute-continuity-uniform}与在一点连续的线性算子的一致连续性的相似性.
\end{remark}

\begin{theorem}[Nikodym]\label{theorem:ch18-nikodym}
令 $(X,\mathcal{M})$ 为可测空间, 而 $\{\nu_{n}\}$ 是 $\mathcal{M}$ 上集合态地收敛到集函数 $\nu$ 的有限测度序列. 假设 $\{\nu_{n}(X)\}$ 是有界的, 则 $\nu$ 是 $\mathcal{M}$ 上的测度.
\end{theorem}

\begin{proof}
对可测集 $E$, 定义
\begin{align}
\mu(E)=\sum_{n=1}^{\infty}\frac{1}{2^{n}}\cdot\nu_{n}(E).\label{align:ch18-nikodym-proof1}
\end{align}
我们把证明 $\mu$ 是 $\mathcal{M}$ 上的有限测度留作练习. 显然, 每个 $\nu_{n}$ 关于 $\mu$ 绝对连续. 结论现在从 Vitali-Hahn-Saks 定理得出.
\end{proof}

\begin{exercise}\label{exercise:ch18-34}
对有限测度空间 $(X,\mathcal{M},\mu)$ 中的可测集 $A$ 和 $B$, 证明:
\begin{align*}
\rho_{\mu}(A,B)=\mu(A)+\mu(B)-2\cdot\mu(A\cap B).
\end{align*}
\end{exercise}

\begin{exercise}\label{exercise:ch18-35}
令 $\{A_{n}\}$ 为关于 Nikodym 度量收敛到可测集 $A_{0}$ 的可测集序列. 证明 $A_{0}\approx\bigcup_{n=1}^{\infty}\left[\bigcap_{k=n}^{\infty}A_{k}\right]$.
\end{exercise}

\begin{exercise}\label{exercise:ch18-36}
令 $(X,\mathcal{M},\mu)$ 为有限测度空间而 $\nu:\mathcal{M}\to[0,\infty)$ 是有限可加集函数, 它具有以下性质: 对每个 $\varepsilon>0$, 存在 $\delta>0$ 使得对可测集 $E$, 若 $\mu(E)<\delta$, 则 $\nu(E)<\varepsilon$. 证明 $\nu$ 是 $\mathcal{M}$ 上的测度.
\end{exercise}

\begin{exercise}\label{exercise:ch18-37}
令 $(X,\mathcal{M})$ 为可测空间而 $\{\nu_{n}\}$ 是 $\mathcal{M}$ 上的有限测度序列, 它在 $\mathcal{M}$ 上集合态地收敛于 $\nu$. 令 $\{E_{k}\}$ 为具有空的交的递减可测集序列. 证明对每个 $\varepsilon>0$, 存在自然数 $k$ 使得对所有 $n$, $\nu_{n}(E_{k})<\varepsilon$.
\end{exercise}

\begin{exercise}\label{exercise:ch18-38}
给出可测空间上递减的测度序列 $\{\mu_{n}\}$, 使得定义为 $\mu(E)=\lim_{n\to\infty}\mu_{n}(E)$ 的集函数 $\mu$ 不是一个测度.
\end{exercise}

\begin{exercise}\label{exercise:ch18-39}
令 $(X,\mathcal{M})$ 为可测空间而 $\{\mu_{n}\}$ 是 $\mathcal{M}$ 上的测度序列, 使得对每个 $E\in\mathcal{M}$, $\mu_{n+1}(E)\geqslant\mu_{n}(E)$. 对每个 $E\in\mathcal{M}$, 定义 $\mu(E)=\lim_{n}\mu_{n}(E)$. 若 $\mu(X)<\infty$, 证明 $\mu$ 是 $\mathcal{M}$ 上的测度.
\end{exercise}

\begin{exercise}\label{exercise:ch18-40}
对带号测度叙述和证明 Vitali-Hahn-Saks 定理.
\end{exercise}

\end{document}